\begin{exercice}\label{exo4-11} (\minsyndical)

Soient $a<b$ deux r\'eels.
Une variable al\'eatoire continue $X : \Omega \rightarrow [a,b]$ suit la {\it loi uniforme} $\mathcal{U}(a,b)$ si sa densit\'e de probabilit\'e est la suivante :

\[
f(x) = \frac1{b-a} \mathds{1}_{[a,b]} (x),
\]

o\`u $\mathds{1}_{[a,b]}$ est la fonction indicatrice de l'intervalle $[a,b]$.

\begin{enumerate}

\item V\'erifier que $f$ est bien une densit\'e de probabilit\'e.

\item Calculer et tracer la fonction de r\'epartition $F_X(t)$ associ\'ee \`a $X$. 

\item Soient $c$ et $d$ tels que $a\leq c < d \leq b$. Calculer $P(c\leq X \leq d)$.\\
Justifier alors le terme de loi uniforme.

\item Calculer l'esp\'erance math\'ematique de la variable $E[X]$.

\item Calculer la variance $V(X)$.

\end{enumerate}

%\corrref{TD1_1}

\end{exercice}
